% --- 03.Mercado.tex ---

\chapter{Mercado Genético del Vacío}

\begin{loreblock}
``La evolución no es un regalo del tiempo — es una inversión de sangre.''
\end{loreblock}

\section*{Evolución por Sacrificio}

Los Biskrorum son \textbf{frágiles en el plano físico} (Resistencia Física baja), pero su \textbf{Resistencia Mágica es abrumadora} y su tecnología de rango es devastadora. Para alcanzar su máximo potencial, deben consumir su propia biomasa — un sacrificio calculado, no emocional.

Al construir tu lista, puedes invertir de \textbf{1 a 3 puntos adicionales} por cada Mutación Genética que equipes en una unidad. Estas modificaciones se registran en la ranura de Equipamiento de la ficha.

\vspace{0.4cm}

\begin{tcolorbox}[enhanced, colback=white, colframe=VoidBlack,
    title={\large \textbf{CATÁLOGO DE MUTACIONES — NIVEL BÁSICO}},
    boxed title style={colback=VoidPurple, coltext=white}]
\begin{tabularx}{\linewidth}{l|c|X}
\rowcolor{BloodPurple!20} \textbf{Gen / Mutación} & \textbf{Pts} & \textbf{Efecto} \\ \hline

\textbf{Glándula de Furia Abisal} & 3 &
\textbf{Sacrificio:} Al inicio de la partida, la unidad pierde \textbf{3 PV de Cuerpo}. \\
& & \textbf{Efecto permanente:} Gana \textbf{+2 al daño de todos sus ataques}. \\ \hline

\textbf{Mutación Óptica} & 2 &
\textbf{Sacrificio:} Pierde \textbf{2 PV de Cabeza} al inicio de la partida. \\
& & \textbf{Efecto permanente:} Gana \textbf{+1 a Range ATK} y puede ignorar penalizadores de cobertura ligera. \\ \hline

\textbf{Espasmo Nervioso} & 1 &
\textbf{Activable (1× por turno):} Sacrifica \textbf{1 PV} para ganar \textbf{+5cm de MOV} inmediato. No requiere PA. \\ \hline

\textbf{Corteza Refractaria} & 2 &
\textbf{Sacrificio:} Pierde \textbf{2 PV de Cuerpo} al inicio de la partida. \\
& & \textbf{Efecto permanente:} Gana \textbf{+2 a Reflejos} contra ataques a distancia. \\ \hline

\textbf{Redundancia Orgánica} & 2 &
Cuando esta unidad llega a 0 PV en \textbf{Brazos o Piernas}, no aplica el efecto de incapacitación inmediata. En cambio, tira 1d6: con \textbf{4+} la extremidad aguanta un turno más antes de fallar. \\ \hline

\textbf{Receptores de Vacío} & 3 &
Esta unidad puede usar \textbf{PE de la Gema} aunque esté fuera del Área de Influencia, siempre que esté a menos de \textbf{30cm del portador}. \\ \hline

\end{tabularx}
\end{tcolorbox}

\vspace{0.4cm}

\begin{tcolorbox}[enhanced, colback=AbyssalBlue!15, colframe=VoidPurple,
    title={\large \textbf{CATÁLOGO DE MUTACIONES — NIVEL AVANZADO}}
    , boxed title style={colback=BloodPurple, coltext=white}]
\begin{tabularx}{\linewidth}{l|c|X}
\rowcolor{VoidPurple!15} \textbf{Gen / Mutación} & \textbf{Pts} & \textbf{Efecto} \\ \hline

\textbf{Tentáculo de Abordaje} & 3 &
Como acción (1 PA), lanza un tentáculo a 8cm. Si impacta (Range ATK vs. Reflejos), arrastra al objetivo \textbf{4cm hacia esta unidad}. El objetivo queda \textbf{Enmarañado} 1 turno. \\ \hline

\textbf{Campo Neuronal} & 4 &
Esta unidad puede transmitir el \textbf{bono de Sincronía del Vacío} a aliados en 20cm aunque no haya un Mago presente, pero solo \textbf{1 vez por ronda} y solo si esta unidad no fue Agotada ese turno. \\ \hline

\textbf{Médula de Titanio} & 2 &
La RD de la zona \textbf{Cuerpo} de esta unidad aumenta en \textbf{+1} permanentemente. Puede apilarse con equipo de armadura. \\ \hline

\textbf{Instinto Suicida} & 1 &
Cuando esta unidad es eliminada, puede gastar \textbf{0 PA} para realizar inmediatamente \textbf{1 ataque Melee} gratuito a cualquier unidad adyacente (aliada o enemiga). Este ataque ignora el estado Agotado. \\ \hline

\end{tabularx}
\end{tcolorbox}

\vspace{0.4cm}

\begin{reglabiskrida}{Regla: Límite de Mutaciones por Unidad}
\begin{itemize}
    \item Cada unidad puede equipar un \textbf{máximo de 3 Mutaciones Genéticas}.
    \item El costo total de las mutaciones de una unidad se suma a su costo base en puntos de lista.
    \item Los sacrificios de PV se aplican al \textbf{inicio de la partida}, antes del primer turno.
    \item Las mutaciones no pueden cambiarse una vez que comienza la partida.
\end{itemize}
\end{reglabiskrida}
